\documentclass[11pt,a4paper]{article}
\usepackage{iftex}
\ifPDFTeX
  \usepackage[T1]{fontenc}
  \usepackage[utf8]{inputenc}
  \usepackage{lmodern}
  \input{glyphtounicode}
  \pdfgentounicode=1
\else
  \usepackage{fontspec}
  % Kerning off: kerned pairs ("A WS", "T ree") split words for ATS text extraction.
  \setmainfont{lmroman10-regular}[Extension=.otf, BoldFont=lmroman10-bold,
    ItalicFont=lmroman10-italic, BoldItalicFont=lmroman10-bolditalic, Kerning=Off]
\fi
\usepackage[margin=0.9in]{geometry}
\usepackage[hidelinks]{hyperref}
\hypersetup{pdftitle={Abhinav Kaushik - Cover Letter - hellofresh}, pdfauthor={Abhinav Kaushik}}

\pagestyle{empty}
\setlength{\parindent}{0pt}
\setlength{\parskip}{10pt}
\raggedright

\begin{document}

{\Large\bfseries Abhinav Kaushik}\\
Senior AI Engineer\\
Bengaluru, India \textbar{} +91 87005 71859 \textbar{} \href{mailto:aabhinav.kaushik@gmail.com}{aabhinav.kaushik@gmail.com} \textbar{} \href{https://www.linkedin.com/in/abhinav-kaushik10/}{LinkedIn}  

September 24, 2026

Hiring Team\\
hellofresh

\textbf{Re: Technical Lead, Menu Personalization ML}

Dear Hiring Team,

HelloFresh needs a unified technical vision and architecture for its end-to-end menu personalization stack. In this role I will define and build a scalable ML pipeline that ingests user interaction data, ranks recipe candidates, and serves low-latency recommendations to millions of customers. My eight years of production AI work, an MSc in Machine Learning and AI, and a track record of shipping complex systems make this credible.

I will first establish the core architecture by extending the LangGraph-based analytics framework I created at Optum, which combined stateful multi-agent reasoning with real-time chart generation and cut manual analysis effort by 98\%. I will replicate that modular, containerised approach for HelloFresh using FastAPI, MLflow and Kubernetes, as I did for a production SLA-breach prediction service at Viavi that reduced downtime by 67\%. To improve relevance I will train a cross-encoder retrieval model, like the one that raised dense-retrieval accuracy by 20\% for AARP survey data, and integrate it into a Spark-driven feature pipeline that I built for a recipe recommendation engine, delivering personalized weekly menus. Observability will be built on the automated evaluation pipelines I delivered at Optum, which cut human-in-the-loop review time and added quality metrics for QA agents.

I move quickly from prototype to production, turning business questions into measurable ML tasks, collaborating with product and engineering, and validating models with automated metrics before release. My experience spans Python, PySpark, Docker, Kubernetes and end-to-end CI/CD pipelines, matching the team's need for reliable, observable services.

My goal is to create a world-class, data-driven recommendation platform that continuously learns from customer behavior and drives higher engagement for HelloFresh. This role offers the perfect arena to apply my expertise in retrieval, ranking and LLM-enhanced reasoning to a consumer-facing product at scale.

My fiancée already lives in Berlin, so this role lets me build my career in the city where we are making our home, and I am committed to settling in Berlin long term. I would welcome the opportunity to discuss how my background can accelerate HelloFresh's personalization ambitions.

Sincerely,\\[4pt]
Abhinav Kaushik

\end{document}